% Generated draft: requires mathematical and source review.
\documentclass[11pt]{article}
\usepackage[margin=25mm]{geometry}
\usepackage{fontspec}
\setmainfont{Latin Modern Roman}
\usepackage{amsmath,amssymb,mathtools,unicode-math}
\setmathfont{Latin Modern Math}
\usepackage{enumitem,needspace,xurl,hyperref}
\setlength{\parindent}{0pt}
\setlength{\parskip}{5pt}
\setlength{\emergencystretch}{3em}
\newcommand{\N}{\mathbb{N}}
\newcommand{\Z}{\mathbb{Z}}
\newcommand{\Q}{\mathbb{Q}}
\newcommand{\R}{\mathbb{R}}
\newcommand{\C}{\mathbb{C}}
\newcommand{\F}{\mathbb{F}}
\newcommand{\A}{\mathbb{A}}
\newcommand{\PP}{\mathbb{P}}
\newcommand{\E}{\mathbb{E}}
\newcommand{\eps}{\varepsilon}
\newcommand{\dd}{\,\mathrm{d}}
\begin{document}
\section*{UnsolvedMath ID: 700801}
\textbf{Source problem code:} catalogue entry
\subsection{Definitions and mathematical statement}
Let \((G_i)_{i\in I}\) be an ascending series of groups, so that \(G_i\leq G_j\) whenever \(i\leq j\), and define its union by
\[
G=\bigcup_{i\in I}G_i.
\]
Assume that, for each \(i\in I\), the group \(G_i\) is a group of Lie type over a field \(K_i\), and that all the groups \(G_i\) have the same Lie rank \(r\).

The question of R. Phillips is whether there must exist a field \(K\) and a group \(H\) of Lie type over \(K\), of rank \(r\), such that
\[
G\cong H.
\]
The fields \(K_i\) are not assumed in the supplied statement to form a chain of subfields, and the group inclusions are not stated to arise from field embeddings.

This formulation retains unresolved terminology from the source: the permitted index sets for the ascending series, the precise class of groups called groups of Lie type, and the convention for rank are not specified. In particular, the source does not explicitly require a common Lie type beyond equality of ranks. The question above is therefore a statement schema subject to clarification of these conventions.

A field is locally finite if every subfield generated by finitely many elements is finite. The supplied source states that the answer is affirmative when all the fields \(K_i\) are locally finite, attributing this case to V. Belyaev, A. Borovik, B. Hartley \& G. Shute, and S. Thomas. This is a supplied qualification; the precise hypotheses and conclusions of the cited results are not included in the record.
\subsection{Short English statement}
Must the union of an ascending series of groups of Lie type, all of the same rank, itself be a group of Lie type of that rank over some field? The source reports an affirmative answer for locally finite fields. The meanings of group of Lie type, rank, and ascending series require clarification before a fully unambiguous formulation is possible.
\subsection{Sources}
\begin{itemize}
\item Local UnsolvedMath registry, supplied complete source record: permanent ID \(700801\), exact source problem code \texttt{catalogue entry}, title \emph{Kourovka Notebook Problem 11.109}. The registry labels the problem \texttt{open}; this label is supplied metadata, not an independently verified assessment.
\item \emph{The Kourovka Notebook (Unsolved Problems in Group Theory)}, original problem 11.109, 11th issue (1990); proposed by R. Phillips. The supplied background identifies the 21st issue edition (2026), \texttt{arXiv:1401.0300v48}. The collection metadata names E. I. Khukhro and V. D. Mazurov as editors.
\item The supplied statement attributes the affirmative locally finite field case to V. Belyaev, A. Borovik, B. Hartley \& G. Shute, and S. Thomas. No individual publication details for these results are supplied.
\end{itemize}
\end{document}
